% GENERATED FILE. Edit canonical lessons, then run npm run generate:books.
% canonical-source-hash: fnv1a64:ec443b9c758bf288
% canonical-lessons: UR-C529-jhamela, UR-C529-makhmasa, UR-C529-kashmakash, UR-C529-pechida, UR-C529-uljhana

\chapter{A Tangle of Troubles}
\label{ch:ur-a-tangle-of-troubles}
\hlchaptermodality{529}

\begin{chapteropening}
\textbf{By the end of this chapter:} \emph{I can talk about awkward choices and muddles.}

Complete the last lesson of chapter 529: I can talk about awkward choices and muddles.
\end{chapteropening}

\section[\texorpdfstring{\ur{جھمیلا} (jhamelā)}{جھمیلا (jhamelā)}]{\ur{جھمیلا} (jhamelā) — a tiresome muddle; a hassle}

\label{lesson:UR-C529-jhamela}



\begin{quote}
Before the new one: say the Urdu for an offering, then the Urdu for to bestow favour on.
\end{quote}

\subsection*{You'll want to know: \ur{جھمیلا}}
\textbf{\ur{جھمیلا}} — \emph{jhamelā} — \textquotedblleft{}a tiresome muddle; a hassle\textquotedblright{}.

\begin{quote}
\ur{اگر} \ur{ٹکٹ} \ur{کھو} \ur{گیا} \ur{تو} \ur{بڑا} \ur{جھمیلا} \ur{ہو} \ur{گا۔}
\end{quote}

\begin{quote}
\emph{agar ṭikaṭ kho gayā to baṛā jhamelā ho gā.}
\end{quote}

\begin{quote}
If the ticket gets lost, it will be a real hassle.
\end{quote}

Masculine.

\begin{cousinweb}
\textbf{\ur{جھمیلا}} (\emph{jhamelā}) is inherited: CDIAL lists the Hindi cousin \emph{jhamelā}, \textquotedblleft{}wrangling\textquotedblright{}, under \emph{*jhamm}, \textquotedblleft{}entanglement\textquotedblright{}, a form it reconstructs (CDIAL 5340).
\end{cousinweb}

\begin{grammarlens}[title={What you've built}]
One more word for saying what happens if.
\end{grammarlens}

\subsection*{Your turn}
\begin{itemize}
\raggedright
  \item \emph{Say it:} \emph{jhamelā}
  \item \emph{Say it:} \emph{jhamelā}, once more
  \item \emph{Say it:} \emph{agar ṭikaṭ kho gayā to baṛā jhamelā ho gā.} — If the ticket gets lost, it will be a real hassle.
  \item \emph{Recall:} say the Urdu for an offering, then the Urdu for to bestow favour on, then say \emph{jhamelā} again
\end{itemize}

\begin{tcolorbox}[breakable,colback=teal!4,colframe=teal!35!black,title={Before you move on}]
What does \ur{جھمیلا} mean? (\textquotedblleft{}A tiresome muddle; a hassle\textquotedblright{}.) Say it once more.
\end{tcolorbox}
\section[\texorpdfstring{\ur{مخمصہ} (makhmasā)}{مخمصہ (makhmasā)}]{\ur{مخمصہ} (makhmasā) — a quandary; a fix}

\label{lesson:UR-C529-makhmasa}



\begin{quote}
Before the new one: say the Urdu for to bestow favour on, then the Urdu for a tiresome muddle.
\end{quote}

\subsection*{You'll want to know: \ur{مخمصہ}}
\textbf{\ur{مخمصہ}} — \emph{makhmasā} — \textquotedblleft{}a quandary; a fix\textquotedblright{}.

\begin{quote}
\ur{اگر} \ur{دونوں} \ur{دوست} \ur{کل} \ur{آئے} \ur{تو} \ur{مخمصہ} \ur{ہو} \ur{گا۔}
\end{quote}

\begin{quote}
\emph{agar donoṅ dost kal āe to makhmasā ho gā.}
\end{quote}

\begin{quote}
If both friends come tomorrow, I'll be in a fix.
\end{quote}

Masculine.

\begin{cousinweb}
\textbf{\ur{مخمصہ}} (\emph{makhmasā}) is usually traced, through Persian, to Arabic \emph{makhmaṣa}, \textquotedblleft{}hunger, a hard plight\textquotedblright{}; the sense \textquotedblleft{}a quandary\textquotedblright{} came later, and the sources we use do not show how it arose.
\end{cousinweb}

\begin{grammarlens}[title={What you've built}]
One more word for saying what happens if.
\end{grammarlens}

\subsection*{Your turn}
\begin{itemize}
\raggedright
  \item \emph{Say it:} \emph{makhmasā}
  \item \emph{Say it:} \emph{makhmasā}, once more
  \item \emph{Say it:} \emph{agar donoṅ dost kal āe to makhmasā ho gā.} — If both friends come tomorrow, I'll be in a fix.
  \item \emph{Recall:} say the Urdu for to bestow favour on, then the Urdu for a tiresome muddle, then say \emph{makhmasā} again
\end{itemize}

\begin{tcolorbox}[breakable,colback=teal!4,colframe=teal!35!black,title={Before you move on}]
What does \ur{مخمصہ} mean? (\textquotedblleft{}A quandary; a fix\textquotedblright{}.) Say it once more.
\end{tcolorbox}
\section[\texorpdfstring{\ur{کشمکش} (kashmakash)}{کشمکش (kashmakash)}]{\ur{کشمکش} (kashmakash) — an inner struggle; a tug-of-war}

\label{lesson:UR-C529-kashmakash}



\begin{quote}
Before the new one: say the Urdu for a tiresome muddle, then the Urdu for a quandary.
\end{quote}

\subsection*{You'll want to know: \ur{کشمکش}}
\textbf{\ur{کشمکش}} — \emph{kashmakash} — \textquotedblleft{}an inner struggle; a tug-of-war\textquotedblright{}.

\begin{quote}
\ur{اگر} \ur{وہ} \ur{گھر} \ur{چھوڑتا} \ur{تو} \ur{دل} \ur{میں} \ur{کشمکش} \ur{ہوتی۔}
\end{quote}

\begin{quote}
\emph{agar voh ghar chhoṛtā to dil meṅ kashmakash hotī.}
\end{quote}

\begin{quote}
If he were leaving home, he would be torn inside.
\end{quote}

Feminine.

\begin{cousinweb}
\textbf{\ur{کشمکش}} (\emph{kashmakash}) is usually traced to Persian \emph{kashmakash}, \textquotedblleft{}pulling to and fro\textquotedblright{}, built on \emph{kash-}, \textquotedblleft{}to pull\textquotedblright{}.
\end{cousinweb}

\begin{grammarlens}[title={What you've built}]
One more word for saying what happens if.
\end{grammarlens}

\subsection*{Your turn}
\begin{itemize}
\raggedright
  \item \emph{Say it:} \emph{kashmakash}
  \item \emph{Say it:} \emph{kashmakash}, once more
  \item \emph{Say it:} \emph{agar voh ghar chhoṛtā to dil meṅ kashmakash hotī.} — If he were leaving home, he would be torn inside.
  \item \emph{Recall:} say the Urdu for a tiresome muddle, then the Urdu for a quandary, then say \emph{kashmakash} again
\end{itemize}

\begin{tcolorbox}[breakable,colback=teal!4,colframe=teal!35!black,title={Before you move on}]
What does \ur{کشمکش} mean? (\textquotedblleft{}An inner struggle; a tug-of-war\textquotedblright{}.) Say it once more.
\end{tcolorbox}
\section[\texorpdfstring{\ur{پیچیدہ} (pechīdā)}{پیچیدہ (pechīdā)}]{\ur{پیچیدہ} (pechīdā) — complicated}

\label{lesson:UR-C529-pechida}



\begin{quote}
Before the new one: say the Urdu for a quandary, then the Urdu for an inner struggle.
\end{quote}

\subsection*{You'll want to know: \ur{پیچیدہ}}
\textbf{\ur{پیچیدہ}} — \emph{pechīdā} — \textquotedblleft{}complicated\textquotedblright{}.

\begin{quote}
\ur{اگر} \ur{سوال} \ur{پیچیدہ} \ur{ہے} \ur{تو} \ur{استاد} \ur{سے} \ur{پوچھیے۔}
\end{quote}

\begin{quote}
\emph{agar savāl pechīdā hai to ustād se pūchhiye.}
\end{quote}

\begin{quote}
If the question is complicated, ask the teacher.
\end{quote}

\begin{cousinweb}
\textbf{\ur{پیچیدہ}} (\emph{pechīdā}) is usually traced to Persian \emph{pēchīda}, \textquotedblleft{}twisted\textquotedblright{}, built on \emph{pēch}, \textquotedblleft{}a twist, a coil\textquotedblright{}.
\end{cousinweb}

\begin{grammarlens}[title={What you've built}]
One more word for saying what happens if.
\end{grammarlens}

\subsection*{Your turn}
\begin{itemize}
\raggedright
  \item \emph{Say it:} \emph{pechīdā}
  \item \emph{Say it:} \emph{pechīdā}, once more
  \item \emph{Say it:} \emph{agar savāl pechīdā hai to ustād se pūchhiye.} — If the question is complicated, ask the teacher.
  \item \emph{Recall:} say the Urdu for a quandary, then the Urdu for an inner struggle, then say \emph{pechīdā} again
\end{itemize}

\begin{tcolorbox}[breakable,colback=teal!4,colframe=teal!35!black,title={Before you move on}]
What does \ur{پیچیدہ} mean? (\textquotedblleft{}Complicated\textquotedblright{}.) Say it once more.
\end{tcolorbox}
\section[\texorpdfstring{\ur{الجھانا} (uljhānā)}{الجھانا (uljhānā)}]{\ur{الجھانا} (uljhānā) — to tangle; to confuse}

\label{lesson:UR-C529-uljhana}



\begin{quote}
Before the new one: say the Urdu for an inner struggle, then the Urdu for complicated.
\end{quote}

\subsection*{You'll want to know: \ur{الجھانا}}
\textbf{\ur{الجھانا}} — \emph{uljhānā} — \textquotedblleft{}to tangle; to confuse\textquotedblright{}.

\begin{quote}
\ur{اگر} \ur{بلی} \ur{نے} \ur{اون} \ur{الجھائی} \ur{تو} \ur{مشکل} \ur{ہو} \ur{گی۔}
\end{quote}

\begin{quote}
\emph{agar billī ne ūn uljhāī to mushkil ho gī.}
\end{quote}

\begin{quote}
If the cat tangles the wool, there will be trouble.
\end{quote}

Past, with \ur{نے} (he or she …): \ur{اس} \ur{نے} \ur{دھاگا} \ur{الجھایا}, \emph{us ne dhāgā uljhāyā}.

\begin{cousinweb}
\textbf{\ur{الجھانا}} (\emph{uljhānā}) is the causative of \emph{ulajhnā}, \textquotedblleft{}to get tangled\textquotedblright{}, which CDIAL lists under Sanskrit \emph{uparundhati}, \textquotedblleft{}obstructs\textquotedblright{} (CDIAL 2221).
\end{cousinweb}

\begin{grammarlens}[title={What you've built}]
One more word for saying what happens if.
\end{grammarlens}

\subsection*{Your turn}
\begin{itemize}
\raggedright
  \item \emph{Say it:} \emph{uljhānā}
  \item \emph{Say it:} \emph{uljhānā}, once more
  \item \emph{Say it:} \emph{agar billī ne ūn uljhāī to mushkil ho gī.} — If the cat tangles the wool, there will be trouble.
  \item \emph{Recall:} say the Urdu for an inner struggle, then the Urdu for complicated, then say \emph{uljhānā} again
\end{itemize}

\begin{tcolorbox}[breakable,colback=teal!4,colframe=teal!35!black,title={Before you move on}]
What does \ur{الجھانا} mean? (\textquotedblleft{}To tangle; to confuse\textquotedblright{}.) Say it once more.
\end{tcolorbox}
